%\documentclass[envcountsect,11pt,trans]{beamer}
\documentclass[11pt,aspectratio=169]{beamer}
\usepackage[english]{babel}
\usepackage[ansinew]{inputenc}
\usepackage{times}
\usepackage{xcolor}
\usepackage{graphicx}
\usepackage{rotating}
\usepackage{bbding,pifont} % two dingbat fonts
\usepackage{amsmath}
\usepackage{amsfonts}
\usepackage{amssymb}
\usepackage{fancybox}
\usepackage{epstopdf}
\usepackage{eso-pic}
%\usepackage[round]{natbib}
\usepackage{ngerman}
\usepackage{tcolorbox}
\usepackage{colortbl}
%\usepackage[authordate,bibencoding=auto,strict,backend=biber]{biblatex-chicago}
%\usepackage{natbib}
%\usepackage[style=bath, backend=biber]{biblatex}
\usepackage{ulem}
\usepackage{tikz}
\usepackage{booktabs}
\usepackage{longtable}
\usepackage{hyperref}
\usepackage{array}
\usepackage{appendixnumberbeamer}
\usepackage[T1]{fontenc}
\usepackage{adjustbox}
\usepackage{listings}
\usepackage{uarial}
\usepackage{fancyvrb}
\usepackage{subfigure}
\renewcommand{\familydefault}{\sfdefault}

\newcommand*\circled[1]{\tikz[baseline=(char.base)]{
            \node[shape=circle,draw,inner sep=2pt] (char) {#1};}}
\newcommand{\E}{\operatorname{E}}
\newcommand{\Var}{\operatorname{Var}}
\newcommand{\overbar}[1]{\mkern 1.5mu\overline{\mkern-1.5mu#1\mkern-1.5mu}\mkern 1.5mu}
\DeclareMathOperator*{\argmin}{arg\,min}
\DeclareMathOperator*{\argmax}{arg\,max}
% \use\smallpackage{pgfpages}
% \pgfpagesuselayout{4 on 1}[a4paper,border shrink=10mm,landscape]

%\xdefinecolor{MyColor}{rgb}{0.14,0.17,0.52}
%\xdefinecolor{MyColor}{rgb}{0.29,0.53,0.82}
%\xdefinecolor{MyBlue}{rgb}{0.14,0.17,0.52}
%\xdefinecolor{MyGreen}{rgb}{0.70,0.78,0.14}
%\xdefinecolor{MyAlert}{rgb}{0.29,0.53,0.82}
%\colorlet{mystructure}{MyColor}
%\usecolortheme[named=mystructure]{structure}
%\setbeamercolor{alerted text}{fg=MyAlert}
\definecolor{Black}{RGB}{0,0,0}
\definecolor{RedGIZ}{RGB}{200,15,15}
\setbeamercolor{title}{fg=black}
\setbeamercolor{section in toc}{fg=black}
\setbeamercolor{subsection in toc}{fg=black}
\setbeamercolor{frametitle}{fg=Black}
\setbeamercolor{tableofcontents}{fg=Black}

%\usetheme{Boadilla}
%\usetheme{Madrid}
%\useoutertheme{infolines}
%\usecolortheme{whale}
%\usecolortheme{lily}
\renewcommand{\inserttitlegraphic}{
\parbox[b][2.5cm][t]{4cm}{\includegraphics[width = 4cm, keepaspectratio]{pictures/GIZ_Logo.png}}  \hspace{0.1cm} \parbox[b][2.5cm][t]{4cm}{\includegraphics[width = 4cm, keepaspectratio]{pictures/IWH_Logo_RGB_DE_Grossformat.png}} \hfill \parbox[b][2.5cm][t]{3.21cm}{\includegraphics[keepaspectratio,width=4.00cm]{pictures/seperationbar.png} \\ \includegraphics[width = 4cm, keepaspectratio]{pictures/BMWI_logo.png}}
}
\defbeamertemplate*{title page}{customized}[1][]
{	
	\vspace{2.75cm}
  \usebeamerfont{title}\begin{center}\textbf{\inserttitle}\end{center}\par
  \usebeamerfont{subtitle}\begin{center}\usebeamercolor[fg]{subtitle}\insertsubtitle\end{center}\par
  \usebeamerfont{author}\insertauthor $\, \vert$ \usebeamerfont{date}\insertdate \par
  \usebeamerfont{institute}\insertinstitute\par
	\vspace{0.25cm}
  \usebeamercolor[fg]{titlegraphic}\inserttitlegraphic


}
\setbeamerfont{subtitle}{size=\normalsize}
\setbeamerfont{author}{size=\tiny}
\setbeamerfont{date}{size=\tiny}
\setbeamerfont{institute}{size=\tiny}

\setbeamercovered{transparent}

\setbeamertemplate{footline}[text line]{%
  \parbox{\linewidth}{\vspace*{-8pt} $\vert$ \hspace{1mm} \today \hspace{1mm} $\vert$ \hfill}}
\setbeamertemplate{navigation symbols}{}

\linespread{1.2}





\mode<presentation>{
    \setbeamertemplate{itemize item}{\color{RedGIZ}$\blacksquare$}
    \setbeamertemplate{itemize subitem}{\color{RedGIZ}$\blacktriangleright$}
		}
\title[DGE--CRED]{Dynamic General Equilibrium Model for Climate Resilient Economic Development}
%\subtitle[]{Training}

\author[Christoph Schult]{Andrej Drygalla, Katja Heinisch and Christoph Schult*} \date[July 2020]{July 2020}
\institute[IWH]{Halle Institute for Economic Research}

 
\logo{\begin{tabular}{c}\includegraphics[keepaspectratio,width=1.00cm]{pictures/seperationbar.png} \\ \includegraphics[keepaspectratio,width=1.00cm]{pictures/GIZ_Logo}\end{tabular}}

\usepackage[style=authoryear,maxbibnames=9,maxcitenames=2,backend=bibtex]{biblatex}
\bibliography{references}

% Settings by yw:
% Outline at beginning of each section:


\AtBeginSection[]{	
	{\setbeamertemplate{footline}{}
	\begin{frame}<beamer>{Outline}
		\tableofcontents[sectionstyle=show/hide, subsectionstyle=show/show/hide, subsubsectionstyle=show/show/hide]
		\setbeamertemplate{footline}{}
		\addtocounter{framenumber}{-1}
	\end{frame}}
}

% Packages:
\usepackage{ragged2e}
\usepackage{upquote}
\usepackage{subfigure}
\usepackage{ragged2e}
\lstset{ 
  backgroundcolor=\color{white},
	breaklines=true,
	basicstyle=\tiny
}
% Section numbering:
\setbeamertemplate{section in toc}[sections numbered]
\setbeamertemplate{subsection in toc}[subsections numbered]
\setbeamertemplate{subsubsection in toc}[subsubsections numbered]
\setbeamertemplate{section in toc}[ball]
\setbeamertemplate{subsection in toc}[ball]
\setbeamertemplate{subsubsection in toc}[ball]
\setbeamercolor{section number projected}{bg={RedGIZ}}
\setbeamercolor{subsection number projected}{bg={RedGIZ}}
\setbeamercolor{subsubsection number projected}{bg={RedGIZ}}
%\usefonttheme{serif}
\begin{document}
%\footnotesize
\usebackgroundtemplate{
\vbox to \paperheight{\vspace{0.1cm}\hbox to \paperwidth{\hfil\includegraphics[width=0.975\paperwidth,height = 0.7\paperheight]{pictures/BackgroundGIZ.jpg}\hfil}\vfil
}}

\begin{frame}<presentation>[noframenumbering,plain]
  \titlepage \\
	{\tiny} \scalebox{.4}{* Research assistance by Yoshiki Wiskamp is greatly acknowledged.}
\end{frame}
\usebackgroundtemplate{
}
{\setbeamertemplate{footline}{}
\begin{frame}<presentation>[noframenumbering]
	\frametitle{Outline}
		 \tableofcontents[hideallsubsections]
\end{frame}
}
%%%%%%%%%%%%%%%%%%%%%%%%%%%%%%%%%%%%%%%%%%%%%%%%%%%%%%%%

\section{Installation of Dynare}

\subsection{Installation of Dynare: Introduction}

\begin{frame}<presentation>
\frametitle{{\thesection.\thesubsection} Installation of Dynare: Introduction}
	\begin{itemize}
		\item Packaged versions of Dynare are available for:
		\begin{itemize}
			\item Windows (7, 8.1, 10)
			\item Several GNU/Linux distributions (Debian, Ubuntu, Linux Mint, Arch Linux)
			\item macOS 10.11 or later
			\item Should work on other systems, but some compilation steps are necessary
		\end{itemize}
	\item In order to run Dynare, you need one of the following:
		\begin{itemize}
			\item MATLAB version 7.9 (R2009b) or above
			\item GNU Octave version 4.2.1 or above, with the statistics package from Octave-Forge
		\end{itemize}
	\item For the DGE CRED model only Dynare 4.6 and the newest Octave version are suitable
	\end{itemize}
\end{frame}


\subsection{Installation of Dynare: Windows}

\begin{frame}<presentation>
\frametitle{{\thesection.\thesubsection} Installation of Dynare: Windows (1)}
	\begin{itemize}
			\item Access the Dynare web page (\texttt{https://www.dynare.org}) and click on ``Download v4.6.1'': 
				\begin{figure}
			\frame{\includegraphics[width=9cm]{pictures/Dynare_installation/windows1}}
		\end{figure}
	\end{itemize}
\end{frame}

\begin{frame}<presentation>
	\frametitle{{\thesection.\thesubsection} Installation of Dynare: Windows (2)}
	\begin{itemize}
		\item The following page will be displayed:
		\begin{figure}
			\frame{\includegraphics[width=9cm]{pictures/Dynare_installation/windows2}}
		\end{figure}
	\item Click on ``Dynare 4.6.1 (exe)'' to download the executable installer.
	\end{itemize}
\end{frame}
\begin{frame}<presentation>
	\frametitle{{\thesection.\thesubsection} Installation of Dynare: Windows (3)}
	\begin{itemize}
		\item Open the downloaded executable installer ``\texttt{dynare-4.6.1.exe}''
		\item Follow the displayed instructions
		\begin{figure}
			\frame{\includegraphics[width=6cm]{pictures/Dynare_installation/windows3}}
		\end{figure}

	\end{itemize}
\end{frame}


\begin{frame}<presentation>
	\frametitle{{\thesection.\thesubsection} Installation of Dynare: Windows (4)}
	\begin{itemize}
		\item Finally, choose a folder in which to install Dynare
		
		\begin{figure}
			\frame{\includegraphics[width=6cm]{pictures/Dynare_installation/windows4}}
		\end{figure}
			\item You will later need to tell MATLAB the path where Dynare is installed
	\end{itemize}
\end{frame}

\subsection{Installation of Dynare: macOS}

\begin{frame}<presentation>
	\frametitle{{\thesection.\thesubsection} Installation of Dynare: macOS (1)}
	\begin{itemize}
		\item Access the Dynare web page (\texttt{https://www.dynare.org}) and click on ``Download v4.6.1'':
		\begin{figure}
			\frame{\includegraphics[width=9cm]{pictures/Dynare_installation/mac1}}
		\end{figure}
	\end{itemize}
\end{frame}
\begin{frame}<presentation>
	\frametitle{{\thesection.\thesubsection} Installation of Dynare: macOS (2)}
	\begin{itemize}
		\item The following page will be displayed:
		\begin{figure}
			\frame{\includegraphics[width=9cm]{pictures/Dynare_installation/mac2}}
		\end{figure}
	\item Click on ``Dynare 4.6.1 (pkg)'' to download the installation package.
	\end{itemize}
\end{frame}
\begin{frame}<presentation>
	\frametitle{{\thesection.\thesubsection} Installation of Dynare: macOS (3)}
	\begin{itemize}
		\item Open the downloaded installation package ``\texttt{dynare-4.6.1.pkg}''
		\item In case that the following security notification is displayed:
		\begin{figure}
			\frame{\includegraphics[width=6cm]{pictures/Dynare_installation/mac3}}
		\end{figure}
		\begin{itemize}
			\item Access: ``System Preferences'' --> ``Security \& Privacy'' --> press: ``Open Anyway''
		\end{itemize}
	\end{itemize}
\end{frame}
\begin{frame}<presentation>
	\frametitle{{\thesection.\thesubsection} Installation of Dynare: macOS (4)}
	\begin{itemize}
		\item The following window should appear:
		\begin{figure}
			\frame{\includegraphics[width=6cm]{pictures/Dynare_installation/mac4}}
		\end{figure}
		\item Follow the displayed instuctions.
	\end{itemize}
\end{frame}
\begin{frame}<presentation>
	\frametitle{{\thesection.\thesubsection} Installation of Dynare: macOS (5)}
	\begin{itemize}
		\item If the installation was successfull the following window will appear:
		\begin{figure}
			\frame{\includegraphics[width=6cm]{pictures/Dynare_installation/mac5}}
		\end{figure}
		\item The Dynare folder should now be listed under ``Applications''.
		\item Dynare is now ready to use and the installation package ``\texttt{dynare-4.6.1.pkg}'' can be deleted.
	\end{itemize}
\end{frame}

\section{Set path to Dynare in Matlab}

\subsection{Set path to Dynare in Matlab: permanently}

\begin{frame}<presentation>
\frametitle{{\thesection.\thesubsection} Set path to Dynare in Matlab: permanently (1)}
	\begin{itemize}
			\item On the MATLAB Home tab, in the Environment section, click on Set Path
				\begin{figure}
			\frame{\includegraphics[width=9cm]{pictures/Dynare_installation/setpath1}}
		\end{figure}
	\end{itemize}
\end{frame}

\begin{frame}<presentation>
\frametitle{{\thesection.\thesubsection} Set path to Dynare in Matlab: permanently (2)}
	\begin{itemize}
			\item Click Add Folder
				\begin{figure}
			\frame{\includegraphics[width=9cm]{pictures/Dynare_installation/setpath2}}
		\end{figure}
	\end{itemize}
\end{frame}

\begin{frame}<presentation>
\frametitle{{\thesection.\thesubsection} Set path to Dynare in Matlab: permanently (3)}
	\begin{itemize}
			\item Select the matlab subdirectory of your Dynare installation
				\begin{figure}
			\frame{\includegraphics[width=9cm]{pictures/Dynare_installation/setpath3}}
		\end{figure}
		\item Then select choose folder
	\end{itemize}
\end{frame}

\begin{frame}<presentation>
\frametitle{{\thesection.\thesubsection} Set path to Dynare in Matlab: permanently (4)}
	\begin{itemize}
			\item Apply the setting by clicking Save button
				\begin{figure}
			\frame{\includegraphics[width=9cm]{pictures/Dynare_installation/setpath4}}
		\end{figure}
	\end{itemize}
\end{frame}

\subsection{Set path to Dynare in Matlab: temporarily}

\begin{frame}<presentation>
\frametitle{{\thesection.\thesubsection} Set path to Dynare in Matlab: temporarily}
	\begin{itemize}
		\item Alternatively, the path to Dynare can be set \textit{temporarily} in Matlab
		\begin{itemize}
			\item This allows to switch between different versions of Dynare
		\end{itemize}
		\item The following command has to be entered in Matlab:
		\begin{itemize}
			\item Windows: \texttt{addpath C:\textbackslash dynare\textbackslash4.6.1\textbackslash matlab}
			\item macOS: \texttt{addpath /Applications/Dynare/4.X.Y/matlab}
		\end{itemize}
	\end{itemize}
\end{frame}


%%%%%%%%%%%%%%%%%%%%%%%%%%%%%%%%%%%%%%%%%%%%%%%%%%%%%%%%

\end{document}